\section{ヒント}\label{節:総合ヒント}

以下では,計算の結論ではなく,最初の変形や見るべき条件を示す.解き直すときには,その変形を選んだ理由も言葉にしてみよう.

\begin{enumerate}[label=\textbf{\thechapter.\arabic*},leftmargin=*]
 \item 分子を$1+x^2$で割ると,余りの項が分母の微分と結び付く.
 \item 根号の中を$u=1+x^2$と置く.\,$x\dd x$をそのまま$\dd u$で表す.
 \item $\log x$を微分する側に選ぶと,残る積分は整式になる.
 \item $G(u)=\int_0^u\sqrt{1+t^2}\dd t$と置いて,\,$F(x)=G(x^2)-G(x)$と書く.
 \item $x=-t$で書き直した積分と元の積分を足す.\,$1+x^2$は偶関数である.
 \item 交点を求めた後で,上下関係と,回転軸からの二つの距離を調べる.
 \item 端点を微分すると$H'(a)$が得られる.値の計算には$y=a-\frac12$と置き,積分区間の中心から測ると整理しやすい.
 \item $u=1+x^2$と置けば,\,$x^3\dd x=\frac12(u-1)\dd u$になる.その後で部分積分する.
 \item 一項を$\frac1n$と$\frac{k}{n}$だけを使って表す.積分に直したら,分子を分母と定数の差にする.
 \item 区間$[x-1,x]$は常に$0$を含む.回転後の半径は,両端の軸からの距離の大きい方である.
 \item $K_n+K_{n-1}$では分母が消える.極限には$0\leqq K_n\leqq\int_0^1x^n\dd x$を使う.
 \item $0$付近では$x^\alpha$,無限大では$x^{\alpha-2}$と比較し,二つの条件を同時に満たす範囲を求める.
 \item 積分を$x\int_0^xf(t)\dd t-\int_0^xtf(t)\dd t$に分けて微分する.\,$f''=f$が得られたら,\,$\e^{-x}(f+f')$の微分を考える.
 \item まず二つの制約を満たす一次関数$\ell(x)$を求める.\,$g=f-\ell$と置いたとき,\,$\int_0^1\ell g\dd x$がどうなるかを調べる.
 \item $x'(t)=2\sin^2t$を使う.弧長では$\sqrt{x'(t)^2+y'(t)^2}$を整理してから積分する.
 \item (1) $\sin u$の六次テイラー多項式では,偶数次の係数が$0$になる.七階微分の符号と大きさを調べる.(2) $u=x^2$と置く.収束には部分積分,絶対値の発散には長さ$\pi$の区間ごとの下からの評価を使う.
 \item $x=\sin t$によって$C_n$を$\cos$の奇数べきの積分にする.隣り合うウォリス積分の比と積を組み合わせる.
 \item 指数関数を含む重みは,\,$x$を$-x$にすると相補的になる.その後,\,$x=2t$,さらに$u=\sin t$と置く.残る積分の$(1-u^2)$の指数は整数なので,部分積分で一つずつ下げられる.

 \item 分子と分母をそれぞれ平方に直し,\,$b_{n+1}=b_n^2$を確かめる.
 \item 弦と放物線の差は$(x-a)(b-x)$となる.三角形は,弦の中点を通る縦の線分で二つに分ける.
 \item 二項定理で展開し,端点で残る導関数を調べる.部分積分では$\e^x$を積分し続ける.
 \item $\frac1{m+k}$を$x^{m+k-1}$の積分で表し,二項定理を使う.
 \item $x^2-(ax+b)=p(x)+(1-a)x-\frac16-b$と書き,平方を展開する.
 \item (1) $r=1$を境に,四分円が正方形からはみ出す.(2) 下側の三角形で$y=xt$と置く.内側の積分では$x$を固定する.
\end{enumerate}

\smallskip\noindent 解答は\pageref{解答:第13章}ページ.

